\subsection{Survey Design}

The survey is designed to collect data on students' perceptions of the learner-facing dashboard in a Jupyter-based programming environment.

\subsubsection{Data Collection}

Before the data collection could start, the study was reviewed and approved by the Ethics Committee of the University of Duisburg-Essen. The ethics approval ID is 2601CMPL2910.
The data were collected via an online survey using LimeSurvey, in which participants could take part only once (see Appendix~\ref{sec:appendix_questionnaires}).
This standardized survey design ensured that all participants answered the questions in the same order in order to reduce potential confounding variables that might arise from differing procedures. In this way, a high level of replicability of the study could be ensured.

At the beginning of the study, participants were informed about the purpose of the study and the associated ethical aspects. It was made clear that participation was voluntary and that participants could withdraw from the study at any time without providing reasons or facing any consequences. Prior to starting the survey, participants were therefore required to give their informed consent. The questionnaire then began with initial questions regarding their previous experience with Jupyter Notebook and Python for comparison purposes. After that, the students began testing the Jupyter Dashboard while solving Python exercises.
Subsequently, questionnaires measuring usability and usefulness of the Jupyter Dashboard and a person's subjective mental workload by assessing six key dimensions such as mental demand, physical demand, temporal demand, performance, effort, and frustration level were administered.
This was followed by the collection of sociodemographic information, including age, gender, course of study and what semester the participants are in, in order to better describe and analyze the sample. A control question was added at the end to check whether the subjects had completed the questionnaire honestly and carefully. Finally, participants were informed about the further course of the study and thanked for their participation.

\subsubsection{Measurements}

The usability of the dashboard was assessed using the System Usability Scale (SUS)~\cite{germanupa_sus}. The SUS is a standardized and widely used instrument for evaluating the perceived usability of interactive systems. It consists of 10 items that capture users' subjective assessment of system usability, including aspects such as complexity, ease of use, and confidence in using the system. The items are rated on a 5-point Likert scale ranging from 1 (``strongly disagree'') to 5 (``strongly agree''). The scale includes five positively worded and five negatively worded items, requiring reverse coding of the negatively formulated statements prior to analysis. A sample item is: ``I thought the system was easy to use.'' The SUS score is calculated by converting item responses according to the standard scoring procedure and multiplying the sum by 2.5, resulting in a total score ranging from 0 to 100, with higher values indicating higher perceived usability. A score above 68 is commonly interpreted as indicating above-average usability. The SUS was chosen because it has demonstrated good reliability and validity across a wide range of domains and technologies with a Cronbach's Alpha $=.95$~\cite{chu2019usability}.

Perceived usefulness and acceptance of the dashboard were measured using constructs derived from the Technology Acceptance Model (TAM)~\cite{davis1989technology}. According to the model, individuals' intention to use a system is primarily determined by their perceived usefulness and perceived ease of use. Perceived usefulness refers to the extent to which a person believes that using a particular system enhances their performance, whereas perceived ease of use reflects the degree to which a person believes that using the system is free of effort. All items were rated on a 7-point Likert scale ranging from 1 (``extremely unlikely'') to 7 (``extremely likely''). Example items include statements such as ``Using the dashboard improves my learning performance'' and ``I find the dashboard easy to use.'' For each construct, mean scores were calculated by averaging the corresponding items, with higher values indicating stronger agreement. The TAM has been widely validated across technological and educational contexts and typically demonstrates good internal consistency for both dimensions with a Cronbach's Alpha $=.89$~\cite{masrom2007technology}.

Perceived workload associated with using the dashboard was assessed using the NASA Task Load Index~\cite{hart1988development}.
The NASA-TLX is a widely used multidimensional instrument for measuring subjective workload in human-computer interaction and learning environments. It captures the overall perceived task load across six dimensions: mental demand, physical demand, temporal demand, performance, effort, and frustration~\cite{hart1988development}.
